\chapter{QS-Maßnahmen}
Im Folgenden werden die durchgeführten Maßnahmen zur Qualitätssicherung (QS) der Studienarbeit beschrieben.

Farberklärung:
\begin{itemize}
    \item {\textcolor{yellow!80!brown}{QS-Maßnahme in Arbeit}}
    \item {\textcolor{ForestGreen}{QS-Maßnahme erfolgreich und abgeschlossen}}
    \item {\textcolor{red!80!black}{QS-Maßnahme nicht erfolgreich und abgeschlossen}}
\end{itemize}
\section{Besprechung mit Betreuer}
Regelmäßige Besprechungen mit dem Betreuer der Studienarbeit wurden durchgeführt, um den Fortschritt zu überprüfen, Feedback zu erhalten und sicherzustellen, dass die Arbeit den Anforderungen entspricht. Diese Besprechungen fanden in Form von Online-Meetings und persönlichen Treffen statt.
Bisherige Treffen: 
\begin{itemize}
    \item 28.07.2025: Erstes Treffen zur Besprechung des Themas und der Zielsetzung der Studienarbeit.
    \item 09.10.2025: Treffen an der DHBW zur Besprechung des Projektplans und der ersten Schritte. Übergabe der bestellten Hardware.
\end{itemize}
Teilnehmende: Christoph Niederer, Luca Riebel, Silvan Müller.

\textcolor{yellow!80!brown}{\textbf{QS-Maßnahme in Arbeit}}
\section{Testen der Hardware}
Während der Erstellung des ersten Konzeptes (05.11.2025) wurde der Raspberry Pi getestet, indem der erste Software-Prototyp darauf ausgeführt wurde. Dabei wurde festgestellt, dass der Raspberry Pi einen defekt aufweist, da er die SD Karte nicht auslesen kann. Dies wurde dem Betreuer mitgeteilt, um eine Lösung zu finden.

Durch diese frühzeitige QS-Maßnahme, das testen der Hardware, konnte ein potenzielles Problem erkannt und adressiert werden, bevor es den Fortschritt der Studienarbeit erheblich beeinträchtigt hätte.

Teilnehmende: Christoph Niederer, Luca Riebel.

\textcolor{yellow!80!brown}{\textbf{QS-Maßnahmen in Arbeit}}

% \footnotetext{\textcolor{yellow!80!brown}{QS-Maßnahme in Arbeit}}
% \footnotetext{\textcolor{ForestGreen}{QS-Maßnahme erflogreich und abgeschlossen}}
% \footnotetext{\textcolor{red!80!black}{QS-Maßnahme nicht erflogreich und abgeschlossen}}